\section{Methods}
\label{sec:methods}

\subsection{Study Design and Data Source}

We conducted a retrospective multicenter study using electronic health
record data from six Israeli general hospitals, drawn from Kineret
\citep{rappoport2025kineret}, the
data warehouse of Israel's largest hospital network, represented in the
Observational Medical Outcomes Partnership (OMOP) common data model. Data
were extracted in two batches corresponding to the periods before and after
the national implementation of the inpatient hyperglycemia-treatment
recommendations. The analysis window spanned August 2022 through December
2025.

The assessed recommendations are those of the national position paper on
the treatment of diabetes in hospitalization in a general hospital
\citep{hochberg2023inpatient}, issued in September 2023 by the Israel
Endocrine Society, the Israel Diabetes Association, and the National
Diabetes Council. The pre-intervention period therefore ends immediately
before its publication and the post-intervention period begins a year after
it, with the intervening months excluded to allow for staggered adoption.

The study received Helsinki committee approval covering up to 30,000 adult
patients admitted between 2022 and 2025, defined by the cohort criteria set
out in Section~\ref{sec:cohort-definition}.

\subsection{From OMOP Tables to a Unified Event Dataset}

The source OMOP tables were reorganized into five unified, columnar
tables---visits, visit types, clinical events, drug exposures, and
measurements---without semantic transformation. Measurement units were
normalized at this stage, but source concept identifiers and names were
preserved verbatim, so that all terminology decisions are made explicitly at
the abstraction stage rather than being embedded in the extract.

Kineret was not mapped to an international coding system at the time of this
work, so a dataset-specific concept index---a many-to-one mapping from local
measurement, medication, and observation codes onto the concepts named in
the knowledge base---had to be constructed before abstraction could proceed.
The index is what makes the remainder of the pipeline dataset-agnostic: once
it exists, the same knowledge base and engine run unchanged on any source
within the same clinical scope. Curation was performed jointly by a clinical
domain expert and a data scientist, with candidate codes surfaced by
descriptive analytics over the extracted concept dimensions and adjudicated
by the domain expert. The curated indices, rather than the exploratory
analytics used to generate the candidate lists, are the reproducible
artifact of this step, and are provided as supplementary material.
Further detail on the extraction and on the derived concepts is given in
Appendix~\ref{app:data-prep}.

\subsection{Cohort Definition}
\label{sec:cohort-definition}

An admission was considered eligible if it satisfied all of the following
conditions. The admission was an inpatient visit with at least one
measurement recorded between 18:00 and 07:00, used as evidence of a genuine
overnight stay; its length of stay was between 48 hours and 14 days; the
patient was at least 18 years old at admission; the admission did not take
place in an excluded department; the visit identifier was not associated
with more than one patient identifier, which was treated as a data-integrity
failure; the patient did not die within 48 hours of admission; and the
admission met at least one criterion for diabetes or hyperglycemia.

Exclusions were applied through curated indices rather than inline rules, so
that each exclusion is recorded once alongside the concept vocabulary and
applied uniformly across the pipeline. Pediatric and obstetric departments
were excluded, which together with the adult age criterion also removes
pregnancy-related admissions, including gestational diabetes.

Diabetes or hyperglycemia was established by at least one of three
non-mutually-exclusive criteria: (i) a documented diabetes diagnosis,
identified through a manually curated list of concept identifiers;
(ii) recurrent hyperglycemia, defined as two glucose measurements
$\geq180$ mg/dL separated by at least 2 and at most 24 hours; or
(iii) extreme hyperglycemia, defined as any glucose measurement
$\geq250$ mg/dL. Hypoglycemia was deliberately not used as an entry
criterion: the cohort is defined by hyperglycemic burden, and hypoglycemia
is treated throughout as an outcome rather than as a route into the study
population.

For the pre--post comparison, admissions between August 2022 and July 2023
constituted the pre-intervention period and admissions between August 2024
and December 2025 the post-intervention period. The intervening months were
excluded from the primary comparison to allow for staggered adoption.

Thirty-day mortality was defined from the linked death records as death
occurring between 14 days before and 30 days after the recorded discharge
datetime. The backward tolerance accommodates admissions in which the death
record precedes the administrative discharge timestamp, so that in-hospital
deaths are captured.

\subsection{Constructing the Abstraction Input}

Two inputs were derived from the unified event dataset.

The \emph{temporal input} is a long-format table in which each row records a
single time-stamped clinical fact: admission identifier, concept name, start
and end datetime, and value. A small number of concepts were derived rather
than extracted directly, and structural tokens were inserted to make the
admission timeline explicit to the abstraction engine---an \textsc{admission}
token at the visit start, a terminal \textsc{release} or \textsc{death}
token at its end, and recurring \textsc{meal} tokens at fixed daily times,
which supply the anchors required by recommendations defined relative to
meals. Meal times are not recorded in the source data and were therefore
assumed to occur at fixed daily times---breakfast, lunch, dinner, and a
night snack---following standard hospital practice. Selected concepts
recorded before admission, such as chronic
diagnoses and home medications, were propagated into the admission window
using per-concept relevancy lookback periods
(Appendix~\ref{app:data-prep}).

The \emph{context input} is a static table with one row per admission,
holding baseline characteristics ascertained from the first 48 hours:
demographics, chronic condition flags, and indicators distinguishing
medications taken at home from those administered in hospital. From these
flags a comorbidity score was computed, following the structure of the
Charlson Comorbidity Index and referred to throughout as a modified Charlson
comorbidity index (CCI). The modification was necessary because the source
dataset was not mapped to a standard vocabulary, so condition categories
were available only to the extent that the corresponding codes could be
identified during curation; the weighting was adapted to the categories the
extract supports, and is given in Appendix~\ref{app:data-prep}.

\subsection{Guideline Adherence Assessment}

Guideline adherence was assessed by knowledge-based temporal abstraction
\citep{shahar1997framework}, which converts time-stamped clinical data into
named interval concepts---states, trends, events, and contexts---and then
detects, over those intervals, the temporal pattern that correct execution
of each recommendation would have left in the record
\citep{ben2024supporting}. Abstraction and
assessment are performed by the Temporal Mediator in a single pass, using
our reimplementation of it \citep{mediator-github}. All domain knowledge is
declared externally as XML units validated against a shared schema, so that
both the clinical semantics and the guideline requirements are stated once,
outside the analysis code.

Each recommendation is formalized as a quality-assessment temporal pattern:
an \emph{anchor} identifying the clinical situation in which the
recommendation applies, an \emph{action} that satisfies it, a constraint
type, and a membership function grading partial compliance. Applicability is
enforced by context---an anchor falling outside the context in which the
recommendation holds generates no obligation and is not scored---so that
an insulin dose after a meal, for example, is required only
where glucose levels are elevated. Recurring recommendations such as routine glucose
monitoring are instead evaluated by counting qualifying events within
fixed-length windows across the admission.

Compliance is graded rather than binary. Each constraint carries a
trapezoidal membership function mapping the observed timing, frequency, or
value to $[0,1]$: if correction insulin is expected within 4 hours of a
hyperglycemic reading and considered missed after 12, insulin given at 8
hours scores $0.5$ rather than counting as an outright failure. Three
constraint types are used, corresponding to the compliance aspects reported
in the Results---\emph{time}, the interval between anchor and action;
\emph{value}, the magnitude of the recorded value, most often a dose; and
\emph{cyclic}, the number of qualifying events within a window. Aspect
scores are computed separately and averaged at analysis, and trapezoid
boundaries may be rescaled per patient, which is how weight-normalized
dosing is expressed.

Two conventions govern unmet obligations. Where the recommendation requires
an action, an unanswered anchor scores 0; where it requires an action to be
withheld, the absence of a response is itself compliant and scores 1. A
recommendation that never became applicable yields no output rather than a
zero, so inapplicable recommendations do not depress the score.

Discretization thresholds, relevance durations, and trend limits were
acquired in two passes. One-dimensional $k$-means and kernel density
estimation over each concept's value distribution proposed candidate
boundaries; these were then reconciled against published reference ranges
and the definitions in the national guideline, which were authoritative,
with the data-driven candidates used only to resolve cases where the
literature gives a range rather than a value. Concepts with too many
distinct value combinations for in-engine discretization, such as
corticosteroid preparations at varying doses, were pre-discretized during
data preparation.

The knowledge base was authored by the first author and developed under the
guidance of two co-authors, one specializing in medical informatics and
temporal reasoning and one in diabetes care, who reviewed the formalized
recommendations for clinical fidelity against the source position paper
\citep{hochberg2023inpatient}. Because it names clinical concepts
rather than dataset columns, it is transferable to any electronic health
record within the same clinical scope once a concept index maps that
source's codes onto its vocabulary. Scoring is retrospective and offline:
the procedure issues no recommendation and returns nothing to the treating
team. The engine emits one compliance score per admission, pattern,
compliance aspect, and evaluation time, and these scores are the exposure
variable for all subsequent analyses.

\subsection{Outcome Definitions}

Ten complications were evaluated: 30-day mortality, hyperglycemia, severe
hyperglycemia, hypoglycemia, severe hypoglycemia, hyperosmolality,
ketoacidosis, infection, Inpatient acute kidney injury (AKI), and inpatient myocardial infarction (MI).
All ten are emitted as abstractions by the same engine, so that outcome
ascertainment and adherence assessment draw on a single, explicit knowledge
base. They differ only in where the underlying logic is evaluated. The
glycemic outcomes and acute kidney injury are derived by the engine directly
from measurements: hyperglycemia and hypoglycemia from glucose values, using
a parameterized concept that distinguishes a sustained state from an
isolated reading, and acute kidney injury from serum creatinine and its
ratio to baseline. Ketoacidosis, infection, hyperosmolality, and myocardial
infarction require multi-step logic combining several concept types with
their own temporal windows, and were therefore precomputed during data
preparation and supplied to the engine as primitive concepts, which it
re-emits as complication events. Thirty-day mortality is taken from the
linked death records.

Two outcomes are reported under abbreviated labels and warrant note here.
MI requires evidence of an ischemic
cardiac event together with a troponin level of at least 600, so that coded
cardiac history unaccompanied by biochemical evidence of acute injury does
not qualify. AKI is defined from laboratory
values as stage 2 or above, excluding stage~1. Full operational definitions
of all ten complications are given in Appendix~\ref{app:definitions}, and
the definitions of the assessed guideline recommendations in
Appendix~\ref{app:patterns}.

\subsection{Statistical Analysis}

\subsubsection{Aggregation of adherence scores}

The admission is the unit of analysis throughout. A patient may contribute
several admissions, and patient-level grouping is applied wherever the
statistics require it, as in the clustered standard errors of the mortality
model below. Compliance scores were aggregated in
stages designed so that neither frequently-evaluated patterns nor
frequently-evaluated days dominate the result. Scores were first averaged
within each admission, pattern, compliance aspect, and calendar day; aspects
were then collapsed within each admission, pattern, and day; days were then
collapsed within each admission and pattern; and finally patterns were
averaged with equal weight to give one overall adherence score per
admission. Component-level analyses use the corresponding pattern-level and
aspect-level aggregates.

\subsubsection{Change in adherence following implementation}

Overall and component-level adherence were compared between the pre- and
post-intervention periods using Welch's unequal-variance $t$-test, with each
admission attributed to a period by its admission date. Confidence intervals for the difference in means are Welch
$t$-intervals and are reported unadjusted; $p$-values were adjusted across
the component-level comparisons using the Benjamini--Hochberg false
discovery rate (FDR) procedure. Monthly means are reported with 95\%
confidence intervals derived from the within-month standard error, and the
overall temporal trend is an ordinary least-squares fit to the monthly
means.

\subsubsection{Adherence and 30-day mortality}

Mean whole-admission adherence was compared between admissions ending in
death within 30 days and those not, using Welch's $t$-test. The graded
relationship was then tested by logistic regression on the continuous score,
reported as the odds ratio per $0.1$-point increase; adherence strata appear
in the corresponding figure for display only, with strata under 30
admissions suppressed, and are not the basis of the trend statistic.

The adjusted estimate was a logistic regression on the overall adherence
score with prespecified covariates---age, sex, the modified Charlson
comorbidity index, and hospital as fixed effects---with continuous
predictors standardized, so the odds ratio corresponds to a
one-standard-deviation increase, and standard errors clustered on patient
identity. Both estimates use whole-admission adherence and are reported as
descriptive.

\subsubsection{Component-level models}

To identify which individual recommendations were associated with each
complication, one model was fitted per complication per observation window.
The analysis dataset contains one row per admission and complication,
labelled positive if the complication first occurred at or after the
observation cutoff and negative if it never occurred; admissions in which
the complication occurred \emph{before} the cutoff were excluded, since a
pre-existing condition cannot be attributed to subsequent in-stay care. Both
adherence and baseline covariates were computed only from admission up to
the cutoff, never afterwards. Cutoffs of 24, 48, and 72 hours were used.

Models were fitted by bias-reduced (Firth penalized-likelihood) logistic
regression \citep{firth1993bias}, with adherence and baseline covariates
entered jointly so that each adherence coefficient is estimated holding
measured comorbidity constant. Firth's method was necessary because
unpenalized maximum likelihood failed to converge under separation for most
complications, given the low event counts of several endpoints; the
Jeffreys-prior penalty yields finite estimates and valid Wald inference
without inducing sparsity. Highly correlated predictors were pruned before
fitting and the remainder standardized so coefficients are comparable.
Benjamini--Hochberg FDR correction was applied within each model's own family
of tested predictors, excluding pruned features, which were never tested. A
bootstrap sign-agreement check flagged estimates whose direction was unstable
under resampling. A predictor was considered robust only if it remained
significant with consistent direction across all three windows.

To avoid mechanical circularity, the routine glucose-monitoring component
was excluded from the models for glycemic outcomes, since glucose
measurement is itself the mechanism by which hyperglycemic and hypoglycemic
events are detected.

\subsubsection{Propensity-weighted analysis}

The recurrent components were re-evaluated using inverse probability of
treatment weighting. Adherence was dichotomized at the median of each
outcome's eligible population, so reported coefficients are log-odds ratios
comparing above-median with at-or-below-median adherence, and the
corresponding odds ratio is $\exp(\beta)$. A binary exposure is intrinsic to
this design rather than a simplification: the propensity score is the
probability of belonging to the high-adherence group conditional on baseline
covariates \citep{rosenbaum1983central}, and the standardized mean
differences used to verify balance \citep{austin2009balance} are two-group
contrasts undefined for a continuous exposure. Continuous adherence is
retained in the regression analyses above, making the two designs
complementary: a conditional estimate from regression adjustment and a
balance-verified marginal contrast from weighting.

Propensity scores were estimated by bias-reduced logistic regression on the
baseline covariates, converted to stabilized weights and trimmed at the
first and ninety-ninth percentiles, and weighted logistic regression models
fitted for each outcome. Balance was assessed by the maximum standardized
mean difference after weighting, with $0.10$ required for an estimate to be
interpretable; component--outcome pairs failing this are reported as
non-estimable, since weighting groups that remain dissimilar does not
support the intended contrast.

\subsubsection{Multiplicity}

Benjamini--Hochberg FDR correction \citep{benjamini1995controlling} was
applied within each family of tests:
across component-level pre--post comparisons, across predictors within each
component-level model, and across outcomes within the propensity-weighted
analysis. Confidence intervals are reported unadjusted throughout, so an
interval excluding the null may accompany a non-significant FDR-adjusted
$p$-value.

\subsubsection{Implementation}

All processing and analysis were performed in Python 3.13. The extraction
and cohort-construction pipeline used \texttt{dask} and \texttt{pyarrow}
over Parquet storage, with \texttt{pandas} and \texttt{numpy} for in-memory
operations; the empirical threshold discovery used $k$-means clustering and
kernel density estimation from \texttt{scikit-learn} with peak detection
from \texttt{scipy}.

Guideline adherence was computed using the Mediator abstraction engine
\citep{mediator-github}, which validates the knowledge base against its XML
schema before evaluating the abstractions.

Statistical analysis used \texttt{statsmodels} for the logistic,
propensity-weighted, and Welch comparisons, together with its
Benjamini--Hochberg implementation for multiplicity correction. Bias-reduced
logistic regression was implemented directly, following Firth's iteratively
reweighted least-squares formulation with a Jeffreys-prior penalty, with
standard errors taken from the penalized information matrix; it is released
with the analysis code rather than drawn from a library.

